\section{Numerical stability in floating point}
\label{sec:numerical-stability}

The constructions of this paper are exact: over a finite field, or over $\Q$ in
exact rational arithmetic, the schedules evaluate $P$ without error.  It is
nevertheless natural to ask how they behave in floating-point arithmetic,
where Horner's rule and Estrin's scheme are the standard choices.  We show
that adapted-constant schedules---ours included---do not come with the
coefficientwise backward-stability guarantee of Horner's rule: their
floating-point error bound carries an amplification factor $A$, which we
define and measure below and which is unbounded for our schedules and for
the classical adapted schemes, the floating-point counterpart of the
classical observation that adapted-coefficient methods are prone to lose
significance~\cite[\S4.6.4]{knuth1997seminumerical}.

\subsection{Model and a majorant bound for straight-line programs}

We use the standard model~\cite[Ch.~2--3]{higham2002accuracy}: every
floating-point operation satisfies
$\operatorname{fl}(a\circ b)=(a\circ b)(1+\delta)$ with $\abs{\delta}\le u$,
where $u$ is the unit roundoff, and there is no overflow or underflow.  Set
$\gamma_k=ku/(1-ku)$.  The product lemma states that any product of $k$ factors
$(1+\delta_j)^{\pm1}$ equals $1+\theta_k$ with $\abs{\theta_k}\le\gamma_k$.

\paragraph{What is rounded.}
Preprocessing is exact: the coefficients of $P$ are mapped to the parameters
of the schedule in exact rational arithmetic (or in any precision high enough
to be exact for the purpose).  Evaluation is in floating point: the argument
$x$ arrives as a floating-point number, each parameter is rounded once to the
working precision, $\operatorname{fl}(c_j)=c_j(1+\delta_j)$, and every
arithmetic operation of the schedule rounds, with or without fused
multiply--add.  Every evaluation scheme in this paper is a
\emph{straight-line program} $\mathcal C$: a sequence of wires, each either
the input $x$, a constant $c_j$, or the sum, difference, or product of two
earlier wires.  Integer multiples $k\cdot w$ of a wire count as one product.
Constants that are input data (the coefficients $a_i$ in Horner's rule) are
assumed exactly representable, so that Horner's rule has no preprocessing
rounding at all; the preprocessed constants of our schedules, of
Rabin--Winograd, and of Motzkin--Eve carry the one rounding above.

\begin{definition}[Rounding depth and majorant]
\label{def:rounding-depth}
The \emph{rounding depth} $\rho(w)$ of a wire is defined recursively by
$\rho(x)=0$; $\rho(c_j)=0$ if $c_j$ is exact input data and $1$ otherwise;
$\rho(u\pm v)=\max(\rho(u),\rho(v))+1$; and $\rho(u\cdot v)=\rho(u)+\rho(v)+1$.
The \emph{majorant} $M_{\mathcal C}$ is the program obtained from $\mathcal C$
by replacing every constant $c_j$ by $\abs{c_j}$, every integer multiple
$k\cdot w$ by $\abs{k}\cdot w$, and every subtraction by an addition; it is
evaluated exactly, at $\abs{x}$.
\end{definition}

Note that products \emph{add} rounding depths, so that $\rho(v)$ is exactly the
largest number of rounding factors carried by any tree monomial of $v$; a
max-rule for products would undercount squarings.

\begin{theorem}[Majorant bound]
\label{thm:majorant-bound}
Let $\mathcal C$ compute $P(x)$ exactly in exact arithmetic, with output wire
$w$ satisfying $\rho(w)u<1$, and let $\widehat P(x)$ be the value computed in
floating point.  Then
\[
  \abs{\widehat P(x)-P(x)}\;\le\;\gamma_{\rho(w)}\,M_{\mathcal C}(\abs{x}).
\]
\end{theorem}

\begin{proof}
Expand every wire $v$ as a sum of \emph{tree monomials}: $v=\sum_\pi t_\pi$,
where each $t_\pi$ is a product of leaves ($x$'s and constants) along one
multiplicative path of the expression tree, with the sign carried by the
subtractions.  We claim that the computed value satisfies
$\widehat v=\sum_\pi t_\pi\prod_{j\in J_\pi}(1+\delta_j)$ for multisets
$J_\pi$ of rounding errors with $\abs{J_\pi}\le\rho(v)$ (the same $\delta_j$
may recur when a wire is reused).  For the leaves this is the assumption on
the constants.  For $v=u\pm v'$, the computed
$\widehat v=(\widehat u\pm\widehat v')(1+\delta)$ appends one element to every
$J_\pi$, so the depth increases by one.  For $v=u\cdot v'$, every tree
monomial of $v$ is a product $t_\pi s_\sigma$ of one monomial from each
factor, and $\widehat v=\widehat u\,\widehat v'(1+\delta)$ gives it the
multiset $J_\pi\cup J_\sigma\cup\{\delta\}$ of size at most
$\rho(u)+\rho(v')+1$.  By the product lemma each
$\prod_{j\in J_\pi}(1+\delta_j)=1+\theta_\pi$ with
$\abs{\theta_\pi}\le\gamma_{\rho(w)}$, so
$\abs{\widehat P-P}\le\gamma_{\rho(w)}\sum_\pi\abs{t_\pi}$; and
$\sum_\pi\abs{t_\pi}$ is exactly the value of the majorant program at
$\abs{x}$, because taking absolute values of all leaves and integer multiples
and turning all subtractions into additions removes every cancellation between
tree monomials.  No independence among the $\delta_j$ is used.
\end{proof}

\paragraph{Horner and Estrin.}
For Horner's rule the tree monomials are exactly the terms $a_ix^i$, so
$M(\abs{x})=\sum_i\abs{a_i}\abs{x}^i$ and the bound is the classical
$\gamma_{2n}\sum_i\abs{a_i}\abs{x}^i$ (one fewer factor when $P$ is monic and
the leading multiplication is skipped); the same expansion gives the refined
coefficientwise form $\widehat P=\sum_i a_ix^i(1+\eta_i)$ with
$\abs{\eta_i}\le\gamma_{2i+1}$, and $\gamma_{2n}$ for
$i=n$~\cite[\S5.1]{higham2002accuracy}.  Estrin's scheme on a full tree of
degree $n=2^m-1$ also has $M(\abs{x})=\sum_i\abs{a_i}\abs{x}^i$, and counting
the roundings on the path of $a_ix^i$---$m$ additions, $s_2(i)$ explicit
multiplications, and $i-s_2(i)$ factors inherited from the repeated squarings
$x^2,x^4,\dots$, where $s_2(i)$ is the number of ones in the binary expansion
of $i$ (here the additive rule for products matters: the computed
$x^4=x^4(1+\delta_1)^2(1+\delta_2)$ carries three factors)---gives
$\abs{\eta_i}\le\gamma_{i+m}$, a uniform bound $\gamma_{n+m}$ that is
\emph{smaller} than Horner's $\gamma_{2n}$ (again one fewer when monic).  With
a fused multiply--add at every node the product roundings disappear and both
schemes have the uniform bound $\gamma_n$, Horner retaining slightly sharper
coefficientwise constants $\gamma_{\min(i+1,n)}$ versus
$\gamma_{i+m-s_2(i)}$.  Both schemes are therefore \emph{coefficientwise
backward stable}: the computed value is the exact value of a polynomial with
relatively perturbed coefficients, and the relative forward error is at most
$\gamma_\rho\,\kappa(P,x)$ for the condition number
$\kappa(P,x)=\sum_i\abs{a_i}\abs{x}^i/\abs{P(x)}$.  The comparison is of
practical interest in production libraries: the discussion accompanying the
Estrin implementation in Boost.Math~\cite{boostmath932} conjectures bounds of
this shape, and the display above makes them precise---Estrin's tree does not
cost accuracy; its uniform constant $\gamma_{n+m}$ is in fact the smaller
one.

\subsection{Adapted-coefficient schemes and amplification}

For a schedule whose constants are not the coefficients themselves, the tree
monomials are products of $x$ with preprocessed constants, and they cancel: the
majorant can be much larger than $\sum_i\abs{a_i}\abs{x}^i$.  Define the
\emph{schedule amplification}
\[
  A_{\mathcal C}(P,x)\;=\;\frac{M_{\mathcal C}(\abs{x})}{\sum_{i}\abs{a_i}\abs{x}^i}.
\]
Theorem~\ref{thm:majorant-bound} then gives the relative forward bound
$\abs{\widehat P-P}/\abs{P}\le\gamma_{\rho}\,A_{\mathcal C}(P,x)\,\kappa(P,x)$,
and $A_{\mathcal C}$ is precisely the price paid for the multiplications
saved; for Horner and Estrin $A\equiv1$.  For an adapted scheme it is
unbounded as soon as a low coefficient is obtained by cancellation.  Write
$c=(c_1,\dots,c_m)$ for the constants of $\mathcal C$, regarded as
indeterminates, so that the coefficients $a_j=a_j(c)$ of the computed
polynomial $P[c]$ are polynomials in $c$, and call $c$ \emph{admissible} if
the scheme produces it: for our schedules $c=c(\alpha)$ for a key vector
$\alpha\in K^n$, each constant being a key or a signed integer combination
of a few keys; for Rabin--Winograd and Motzkin--Eve, $c$ ranges over the
parameter vectors of the scheme.

\begin{proposition}[Unit pivots make the amplification unbounded]
\label{prop:unit-pivot-unstable}
Suppose that some coefficient has an additive unit pivot,
$a_d=c_i+f(c_{\mathrm{rest}})$ with $f$ independent of $c_i$ (a pivot of
slope $\lambda\ne0$, $a_d=\lambda c_i+f$, is handled identically with
$\lambda c_i$ in place of $c_i$).
\begin{enumerate}[label=(\roman*)]
\item\label{it:pivot-lower} For every $c$ and every $x\ne0$,
  $M_{\mathcal C}(\abs{x})\ge\abs{c_i}\abs{x}^d$, hence
  \[
    A_{\mathcal C}(P[c],x)\;\ge\;\frac{\abs{c_i}\abs{x}^d}{\sum_j\abs{a_j}\abs{x}^j}.
  \]
\item\label{it:pivot-sup} If some admissible $c$ has
  $a_0(c)=\dots=a_d(c)=0$ but a nonzero tree monomial of degree $d$---for
  instance $c_i\ne0$, i.e.\ $f(c_{\mathrm{rest}})\ne0$---then
  $A_{\mathcal C}(P[c],x)\to\infty$ as $x\to0$, and in particular
  $\sup_{P,x}A_{\mathcal C}(P,x)=\infty$.  This is the case whenever $f$
  does not vanish identically on the admissible $c$ (for $c=c(\alpha)$:
  $f(c(\alpha))\not\equiv0$ as a polynomial in $\alpha$) and either $d=0$
  and $c_i$ can be varied independently of
  the other constants, or the coefficient map is onto the monic polynomials
  of degree $n$ and $f(c_{\mathrm{rest}})$ depends on $c$ only through
  $a_{d+1},\dots,a_{n-1}$, as it does at row $d$ of a descending triangular
  decoder.
\end{enumerate}
\end{proposition}

\begin{proof}
\ref{it:pivot-lower}: Group the tree monomials of $P[c]$ by their degree in
$x$.  Those of degree $d$ sum to $a_d(c)\,x^d$, identically in $c$; the ones
among them that contain no leaf $c_i$ sum to $g(c_{\mathrm{rest}})\,x^d$ for
some polynomial $g$, and the ones that do contain $c_i$ sum to
$(a_d-g)\,x^d$.  Setting $c_i=0$ kills every monomial of the second kind, so
$g=a_d|_{c_i=0}=f$, and the tree monomials of degree $d$ containing $c_i$
sum to $c_i\,x^d$.  The majorant is the sum of the absolute values of all
tree monomials (proof of Theorem~\ref{thm:majorant-bound}), so
$M_{\mathcal C}(\abs{x})\ge\abs{c_i}\abs{x}^d$.

\ref{it:pivot-sup}: Write $M_{\mathcal C}(t)=\sum_e m_e t^e$, where
$m_e\ge0$ is the sum of the absolute values of the tree monomials of degree
$e$; the hypothesis says $m_d>0$ (and $m_d\ge\abs{c_i}$ by
\ref{it:pivot-lower}).  At such a $c$ the polynomial is
$P[c]=x^n+\sum_{d<j<n}a_jx^j$, so for $0<\abs{x}\le1$
\[
  A_{\mathcal C}(P[c],x)\;\ge\;
  \frac{m_d\abs{x}^d}{\abs{x}^n+\sum_{d<j<n}\abs{a_j}\abs{x}^j}
  \;\ge\;\frac{m_d}{\abs{x}\bigl(1+\sum_{d<j<n}\abs{a_j}\bigr)}
  \;\longrightarrow\;\infty\qquad(x\to0).
\]
For the sufficient conditions, pick an admissible $c^*$ with
$f(c^*_{\mathrm{rest}})\ne0$, which exists by hypothesis.  If $d=0$
and $c_i$ is free, replace $c^*_i$ by $-f(c^*_{\mathrm{rest}})$: then
$a_0=0$ and $c_i\ne0$.  If the coefficient map is onto and $f$ is
determined by the top coefficients, let $c$ be admissible with
$P[c]=x^n+\sum_{d<j<n}a_j(c^*)\,x^j$: then $a_0(c)=\dots=a_d(c)=0$,
$f(c_{\mathrm{rest}})=f(c^*_{\mathrm{rest}})\ne0$, and $a_d(c)=0$ forces
$c_i=-f(c_{\mathrm{rest}})\ne0$.
\end{proof}

\begin{remark}[The pivot alone does not suffice]
\label{rem:unit-pivot-counterexample}
The hypothesis of part~\ref{it:pivot-sup} cannot be dropped: a unit pivot
with $f\not\equiv0$ does not by itself force $\sup A=\infty$.  The chain
$w_1=x+c_1$, $w_2=x+c_2$, $P=x\,w_1w_2+c_0$ computes
$x^3+(c_1+c_2)\,x^2+c_1c_2\,x+c_0$ with two multiplications, and
$a_2=c_1+c_2$ is a unit pivot in $c_1$ with $f=c_2\not\equiv0$; yet, with
$t=\abs{x}$,
\[
\begin{aligned}
  M_{\mathcal C}(t)-\sum_j\abs{a_j}t^j
  &=\bigl(\abs{c_1}+\abs{c_2}-\abs{c_1+c_2}\bigr)\,t^2
  \le2\min(\abs{c_1},\abs{c_2})\,t^2\\
  &\le2\sqrt{\abs{c_1c_2}}\,t^2
  \le t^3+\abs{c_1c_2}\,t ,
\end{aligned}
\]
so $A_{\mathcal C}\le2$ for all real or complex constants and all $x$, with
equality at $c_1=-c_2=1$, $c_0=0$, $\abs{x}=1$.  What fails is exactly the
hypothesis of part~\ref{it:pivot-sup}: $a_1=c_1c_2$ vanishes only when a
factor does, so no admissible $c$ has $a_1=a_2=0$ with $c_1\ne0$.  Over
$\R$ the coefficient map is not onto (its image is $a_2^2\ge4a_1$); over
$\C$ it is onto, but the correction $c_2$ is a root of $z^2-a_2z+a_1$ and
depends on $a_1$ as well, so the decoder is not triangular.  The
unrestricted claim thus fails already for a two-multiplication cubic, the
multiplication count of $P_3$.
\end{remark}

For the schedules of this paper the hypothesis of part~\ref{it:pivot-sup}
is met at $d=0$: the constant term is the last row of the decoder,
$a_0=\alpha_i+f(\alpha_{\mathrm{rest}})$ for one key $\alpha_i$, and the
degree-$0$ tree monomials of the program do not all vanish when $a_0$ does
(for $P_7$ one has $a_0=\alpha_0+\alpha_1\alpha_2+\alpha_1\alpha_4\alpha_5$
and the sum of the absolute values of the degree-$0$ tree monomials is
$\abs{\alpha_0}+\abs{\alpha_1\alpha_2}+\abs{\alpha_1\alpha_4\alpha_5}$, which
at $a_0=0$ is positive as soon as $\alpha_1\alpha_2\ne0$).  We checked this in exact
rational arithmetic, on the programs exactly as the measurements below
evaluate them, at three random key vectors with nonzero entries and $a_0$
forced to zero through $\alpha_i$, for every odd $5\le n\le63$; for even
$n\le62$ the program is $P_n=xP_{n-1}+\alpha_0$, which at $\alpha_0=0$ has
exactly the amplification of $P_{n-1}$; and $P_3$, whose only degree-$0$
tree monomial is its constant term, has a nonzero degree-$1$ tree monomial
at $x^3$ instead.  The polynomial $x^n$ itself is such an instance: our
chains reach it only with nonzero constants (the zero-key skeleton is not
$x^n$; that of $P_7$ is $x^7+2x^6+2x^5+2x^4+x^3+x^2$, and
\Cref{prop:centering} below shows that below $n-1$ multiplications the
skeleton cannot be a pure product tree), and for $3\le n\le31$ we verified
that its majorant keeps a nonzero term of degree $0$ (odd $n\ge5$), $1$
($n=3$ and even $n\ge6$) or $2$ ($n=4$).  Hence at these degrees
$A_{\mathcal C}(x^n,x)\to\infty$ as $x\to0$, and
$\sup_{P,x}A_{\mathcal C}=\infty$.  Rabin--Winograd's scheme at the degrees
$n=2^m-1$ measured here (where the balanced recursion is exact and no
normalization occurs) and Motzkin--Eve's scheme, with its last constant $c_K$
regarded as a free parameter, both have a free leaf constant that enters
$a_0$ with unit slope and a nonzero correction $f$---the constant term of the
innermost monic remainder, respectively $c_K$---so part~\ref{it:pivot-sup}
applies to them as well.
What the proposition
shows is that the forward bound of Theorem~\ref{thm:majorant-bound}---the
guarantee available for such a schedule---is unbounded relative to
$\sum_i\abs{a_i}\abs{x}^i$: none of these schemes comes with a
coefficientwise backward-stability guarantee of the Horner type,
independently of the multiplication count.  An unbounded upper bound is not
by itself a proof of instability; that the loss is real is shown by the
measurements below, which exhibit the errors.  The proposition is a
statement about the supremum (polynomials whose low coefficients vanish),
not a quantitative bound on a box $\abs{a_j}\le1$; the measurements below
show how large $A$ is on such boxes.
The rounding depth, by contrast,
differs between the schemes only by constant factors: every scheme has
$\rho\ge n-1$, because the leading tree monomial $x^n$ alone passes through
$n-1$ rounded products.  Horner attains $2n-1$, Estrin $n+m-1$, and our
schedules measure $3.6n$--$3.8n$ for $n\le31$
(Table~\ref{tab:numstab})---roughly twice Horner.  Logarithmic multiplicative
height shortens the critical path; it does not reduce the number of rounding
factors.  The decisive quantity is therefore $A$, and it depends on where the
data enter.

\paragraph{Two regimes.}
In the \emph{prescribed-keys} regime the schedule's keys---the parameter
vector $\theta$ of the evaluation circuit $E(x;\theta)$ of
\Cref{sec:model}, written $\alpha$ in the constructions---are the data, as in
hashing, where they are random field elements; the polynomial is induced.  The
constants of the schedule are then the keys themselves or signed sums of a few
of them, and $A$ stays moderate at the degrees relevant to hashing (medians
$10^{0.4}$, $10^{1.4}$, $10^{3.4}$, and $10^{6.8}$ at $n=7$, $15$, $31$, and
$63$ in Table~\ref{tab:numstab}); it grows with $n$ because the induced
coefficients are polynomials in the keys whose terms cancel while the
majorant's do not, but the observed double-precision error never exceeds four
units of $u\sum_i\abs{a_i}\abs{x}^i$ up to $n=63$.  In the
\emph{prescribed-coefficients} regime an arbitrary polynomial is given and the
decoder (\Cref{alg:final-decoder}) recovers the keys.  The coefficient map
$\alpha\mapsto(a_0,\dots,a_{n-1})$ is a polynomial automorphism
(\Cref{lem:polynomial-left-inverse-automorphism}), so the decoder is itself a
polynomial map; but it has high degree---as a polynomial in the top
coefficient, degree $96$, $154$, and $1456$ at $n=15$, $19$, and $23$---and
rational coefficients whose denominators are products of the integer pivots of
\Cref{rem:char2} (powers of two at the degrees measured here, $n\le23$, where every internal index is $k=2$).  The degree compounds along three mechanisms of the
construction: monic square roots of top windows (linear growth), subtraction
of squares (doubling), and synthetic division by $x+\beta$ inside the
known-power blocks (geometric in the block length).  For single-digit integer
input coefficients the keys returned at $n=23$ already have thousands of
binary digits, and the low-row keys are exactly the corrections $-f$ of
\Cref{prop:unit-pivot-unstable} that cancel the chain's own tree monomials:
$A$ grows
explosively---its median already exceeds $10^{43}$ at $n=15$---and
floating-point evaluation is useless beyond small degrees, even though the
same schedule evaluates $P$ exactly in rational arithmetic.

\subsection{Measurements}

Table~\ref{tab:numstab} reports, for random monic polynomials of degree
$n\in\{7,15,31\}$ (prescribed coefficients) and $n\in\{7,15,31,63\}$
(prescribed keys) and random dyadic points $x\in[-2,2]$, the rounding depth
$\rho$ of each schedule, the median and maximum of the amplification $A$, and
the observed double-precision forward error in units of
$u\sum_i\abs{a_i}\abs{x}^i$.  The reference values are computed in exact
rational arithmetic.  In the prescribed-coefficients regime the coefficients
are integers in $[-5,5]$; in the prescribed-keys regime the keys are dyadic
rationals $k/16$, $\abs{k}\le16$.  The schedules are generated by the
reference implementation accompanying the paper.  Rabin--Winograd here uses
balanced splitting with one absorbed coefficient per product; its constants
are integer polynomial expressions in the coefficients whose size roughly
doubles with each level of the balanced recursion (for $n=2^m-1$ no division
occurs at all), so its amplification grows, but far more slowly than the
decoder's.  Motzkin--Eve (Theorem~E of
Knuth~\cite[\S4.6.4]{knuth1997seminumerical}) has numerically computed
constants (double-double root finding, printed to 13--17 digits), so its
error column also includes preprocessing error, and degrees at which the
preprocessing failed verification are omitted.  Preprocessed constants are
converted to double with a single correctly rounded conversion, as in the
model; evaluating every product-and-add with a fused multiply--add changes no
entry of the table at the displayed precision, as the majorant analysis
predicts (it lowers $\rho$ by a constant factor and leaves $A$ unchanged); the one deviation is that in the prescribed-keys regime the induced
coefficients are themselves rounded before Horner, Estrin, and Rabin--Winograd
read them, which is negligible at the reported scale.

\begin{table}[htbp]
\centering\footnotesize
\begin{tabular}{@{}llrrrrrr@{}}
\toprule
regime & scheme & $n$ & $\rho$ & $A$ med. & $A$ max & err.\ med. & err.\ max \\
\midrule
prescribed coefficients & Horner & 7 & 13 & $1$ & $1$ & $0$ & $0$ \\
 & Estrin & 7 & 9 & $1$ & $1$ & $0$ & $0$ \\
 & Rabin--Winograd & 7 & 12 & $6.4$ & $184$ & $0$ & $0$ \\
 & Motzkin--Eve & 7 & 27 & $5.5$ & $127$ & $829$ & $2.2\times10^{4}$ \\
 & \textbf{this paper} & 7 & 25 & $10^{12}$ & $10^{27}$ & $0$ & $3.6\times10^{14}$ \\
\addlinespace[2pt]
 & Horner & 15 & 29 & $1$ & $1$ & $0$ & $1.16$ \\
 & Estrin & 15 & 18 & $1$ & $1$ & $0$ & $0.786$ \\
 & Rabin--Winograd & 15 & 22 & $14$ & $10^{3}$ & $0$ & $1.07$ \\
 & Motzkin--Eve & 15 & 59 & $10^{5}$ & $10^{7}$ & $1.8\times10^{3}$ & $1.1\times10^{7}$ \\
 & \textbf{this paper} & 15 & 54 & $10^{43}$ & $10^{71}$ & $2.6\times10^{29}$ & $2.0\times10^{52}$ \\
\addlinespace[2pt]
 & Horner & 31 & 61 & $1$ & $1$ & $0.125$ & $1.37$ \\
 & Estrin & 31 & 35 & $1$ & $1$ & $0$ & $1.13$ \\
 & Rabin--Winograd & 31 & 40 & $10^{4}$ & $10^{8}$ & $0.377$ & $1.1\times10^{3}$ \\
 & Motzkin--Eve & 31 & 123 & $10^{13}$ & $10^{13}$ & $2.5\times10^{10}$ & $1.7\times10^{12}$ \\
 & \textbf{this paper} & 31 & 118 & $10^{4849}$ & $10^{5756}$ & overflow & overflow \\
\addlinespace[2pt]
\midrule
prescribed keys & Horner & 7 & 13 & $1$ & $1$ & $0$ & $0$ \\
 & Estrin & 7 & 9 & $1$ & $1$ & $0$ & $0$ \\
 & Rabin--Winograd & 7 & 12 & $5.7$ & $10^{5}$ & $0$ & $0$ \\
 & Motzkin--Eve & 7 & 27 & $13$ & $379$ & $961$ & $2.9\times10^{4}$ \\
 & \textbf{this paper} & 7 & 25 & $2.3$ & $8.4$ & $0$ & $0$ \\
\addlinespace[2pt]
 & Horner & 15 & 29 & $1$ & $1$ & $0.0236$ & $1.09$ \\
 & Estrin & 15 & 18 & $1$ & $1$ & $0$ & $0.557$ \\
 & Rabin--Winograd & 15 & 22 & $10^{4}$ & $10^{8}$ & $738$ & $3.0\times10^{7}$ \\
 & Motzkin--Eve & 15 & 59 & $10^{5}$ & $10^{7}$ & $3.0\times10^{3}$ & $4.7\times10^{6}$ \\
 & \textbf{this paper} & 15 & 54 & $28$ & $292$ & $0$ & $1.4$ \\
\addlinespace[2pt]
 & Horner & 31 & 61 & $1$ & $1$ & $0.083$ & $2.23$ \\
 & Estrin & 31 & 35 & $1$ & $1$ & $0.188$ & $2.05$ \\
 & Rabin--Winograd & 31 & 40 & $10^{27}$ & $10^{47}$ & $6.4\times10^{25}$ & $1.1\times10^{37}$ \\
 & Motzkin--Eve & 31 & 123 & $10^{10}$ & $10^{12}$ & $3.9\times10^{8}$ & $2.5\times10^{11}$ \\
 & \textbf{this paper} & 31 & 118 & $10^{3}$ & $10^{4}$ & $2.9\times10^{-3}$ & $3.07$ \\
\addlinespace[2pt]
 & Horner & 63 & 125 & $1$ & $1$ & $0.137$ & $4.63$ \\
 & Estrin & 63 & 68 & $1$ & $1$ & $0.121$ & $4.63$ \\
 & Rabin--Winograd & 63 & 74 & $10^{183}$ & $10^{297}$ & $2.2\times10^{115}$ & $5.4\times10^{205}$ \\
 & \textbf{this paper} & 63 & 266 & $10^{7}$ & $10^{10}$ & $9.3\times10^{-5}$ & $3.47$ \\
\addlinespace[2pt]
\bottomrule
\end{tabular}
\caption{Rounding depth $\rho$, schedule amplification $A$, and observed double-precision forward error of the evaluation schemes on random monic polynomials, in two regimes. Errors (err.) are the observed double-precision forward error in units of $u\sum_i|a_i||x|^i$; a coefficientwise backward-stable scheme reports at most about $\rho$. $\rho$ is the rounding depth on generic constants (a constant that happens to vanish drops a rounding). Medians and maxima over 12 polynomials and 3 points each; the Motzkin--Eve rows are over the samples whose numeric preprocessing verified (all $36$ at $n=7$ and $15$; $9$ and $15$ of the $36$ at $n=31$ in the two regimes).}
\label{tab:numstab}
\end{table}


Three features stand out.  First, Horner and Estrin behave as the theory
predicts, with errors of at most a few units.  Second, Rabin--Winograd sits in
between when the coefficients are prescribed, its moderate amplification
matching the growth of its integer constants; but those constants are built
from the coefficients, so in the prescribed-keys regime---where the induced
coefficients have denominators up to $16^n$---it loses significance at
$n=31$ and its errors reach $5.4\times10^{205}$ units at $n=63$, while our schedules, whose constants are the
keys, stay within four units.  Third, in the prescribed-coefficients regime the adapted schemes
lose all significance at moderate degrees---our decoder's exact rational keys
are the most extreme case, with every double-precision evaluation overflowing
by $n=31$---whereas in the prescribed-keys regime the same schedules evaluate
the induced polynomial to a small multiple of $u\sum_i\abs{a_i}\abs{x}^i$ at
the degrees shown.  At $n=7$ the keys of most of our sampled prescribed-coefficient
instances happened to be dyadic rationals of at most 53 bits, so the median
error is zero despite the large latent amplification; the maximum,
$3.6\times10^{14}$ units, comes from instances where they were not, and
exactness is not guaranteed even at that degree.

\subsection{Discussion}

The multiplication count $\floor{n/2}+1$ is achieved by an exact change of
coordinates on the coefficient space, and floating-point evaluation of
arbitrary prescribed coefficients is not among its intended uses; there the
methods of choice remain Horner, Estrin, and their compensated
variants~\cite{higham2002accuracy,yu2025compensated}.  The same phenomenon,
to varying degrees, affects the adapted-coefficient schemes since
Motzkin~\cite{motzkin1955evaluation} and Eve~\cite{eve1964evaluation}, and it
is why correctly rounded libraries treat the evaluation scheme as part of the
polynomial generation~\cite{dinechin2007fast,aanjaneya2023fast}.  Conversely,
whenever the keys are the data, the schedules are accurate in floating point
at hashing-relevant degrees; over finite fields both hashing routes of
\Cref{sec:model} are exact, and the floating-point question arises only over
$\R$, where the natural route is to sample the keys directly.

\paragraph{What helps, and what does not.}
We measured the natural modifications of the construction.  The amplification
is carried by the gadget lanes that square: the $n\equiv3\pmod 4$ pairs and
the $n=15,27,31$ specials, versus the $4k+1$ family
($\log_{10}A$ medians $85$ and $1201$ at $n=19$ and $23$ against $1.5$, $4$,
and $14$ at $n=9$, $13$, and $21$).  Routing $n\equiv3\pmod4$ as
$P=x^2P_{n-2}+a_1x+a_0$ with $P_{n-2}$ in the $4k+1$ family costs one extra
multiplication ($\floor{n/2}+2$, still $\log_2(n+1)-3$ below Rabin--Winograd)
and lowers $\log_{10}A$ from $156$ to $28$ at $n=19$ and from $2131$ to $40$
at $n=23$; double-precision results are then exact at $n=11$ and within
$\sim10^2$ units at $n=15$, but no variant is usable in double precision much
beyond $n\approx15$.  A preprocessing-time guard---the majorant evaluated once
in log scale at the largest $\abs{x}$ of the domain, falling back to Horner
when $\rho uA$ exceeds the tolerance---costs no runtime multiplication and is
the practical safeguard.  Rescaling the argument by a power of two is exact (no rounding is
introduced) and is the most effective single modification, at the cost of two
multiplications by constants, $x\cdot s^{-1}$ and $s^n\cdot(\cdot)$---one if
the latter is merged with the leading-coefficient scaling of a non-monic
input; they cannot be absorbed into the chain, whose wire coefficients must
stay integers: decoding $Q(y)=s^{-n}P(sy)$ and evaluating $s^nQ(x/s)$ with
$s=4$ takes the keys at $n=15$ from $10^{44.7}$ to $10^{0.2}$ and $A$ from
$10^{64.5}$ to $10^{6.5}$.  It cannot be pushed
further, and the reason is structural: an automorphism is nearly linear near
its fixed point, and scaling moves every input toward $x^n$, but our chains
with all keys zero do not compute $x^n$ (the zero-key skeleton of $P_7$ is
$x^7+2x^6+2x^5+2x^4+x^3+x^2$), so the structural tree monomials of degree $j$
are amplified by $s^{\,n-j}$ once the keys are small, and $A$ grows again
like $s^n$.  One might hope to remove the floor by \emph{centering} the
construction so that its zero-key chain is a pure product tree; the following
shows this is impossible at any saving over Horner.

\begin{proposition}[Centering costs all the savings]
\label{prop:centering}
Let a chain with $M$ multiplications have the property that every wire is a
monomial in $x$ when all keys are zero.  Then the Jacobian of its coefficient
map at the zero key has rank at most $M+1$.  In particular a polynomial
automorphism with this property has $M\ge n-1$.
\end{proposition}

\begin{proof}
A key $\alpha$ added to a factor of degree $m$ (all keys zero) contributes to
first order $\alpha\cdot c\,x^{n-m}$, where $c$ counts the paths from that
gate to the output, because along every path the remaining factors are
monomials of total degree $n-m$.  A factor is $x$ or one of the $M-1$
non-output gates, so at most $M$ distinct degrees occur, plus the output key at
degree $0$: the linear part spans at most $M+1$ coefficient rows.  An
automorphism has an invertible Jacobian everywhere, so $n\le M+1$.
\end{proof}

So below $n-1$ multiplications the reference point $x^n$ is necessarily
computed with cancellation, the scaling has an optimum ($s\approx4$ at
$n=15$), and the design quantity that remains is the \emph{degree content} of
the skeleton's cancellations---the lower the degrees of the structural tree
monomials, the earlier the $s^{\,n-j}$ growth bites---together with the size
of the keys of $x^n$ itself (in the $4k+1$ family $x^n$ decodes to keys of size
$10^{3.4}$ at $n=13$, a pure normalization artifact of the monic $T_2$ that a
re-normalized gadget would remove).  Dyadic regauging of the wires, by
contrast, is provably a bit-identical no-op; hybrids that use our gadgets only in
low-degree blocks of a balanced Rabin--Winograd tree cannot save
multiplications, because the paper's advantage over Rabin--Winograd is exactly
the $\log_2(n+1)-2$ shared power-chain products of the global construction;
and compensated (double-double) evaluation helps only if the keys themselves
are stored to the higher precision, at $20$--$200\times$ the cost of Horner.
A random search over $2.6\cdot10^5$ four-multiplication septic circuit
structures (about $10^4$ of them monic of degree seven with a full-rank
coefficient map, $37$ decodable by a causal table) separates the two paradigms
without exception.  Thirty-four of the decodable circuits are polynomial
automorphisms---constant pivots, as in this paper---and every one of them is
unstable, the best with median $A=10^{6.6}$ and maximum $10^{17.5}$ on
coefficients in $[-5,5]$; all $4096$ sign variants of $P_7$ are automorphisms
with $A\approx10^{11}$.  The three remaining circuits are dominant but not
surjective (one rational pivot), and they are exactly the stable ones, with
medians $10^{0.6}$, $10^{1.4}$, and $10^{2.1}$.  The best is
\begin{align}
  g_0&=(a_0-x)\,x, & g_1&=(x-g_0+a_1)(x+g_0+a_2),\\
  g_2&=(g_0+a_3)(x+g_1+a_4), & g_3&=x\,(x+g_0-g_1+g_2+a_5),
\end{align}
$P=g_0-g_2+g_3+a_6$: four constant pivots, then a division by a cubic
$\pi(a_4,a_5,a_6)$ in the top coefficients; polynomials with $\pi=0$ are not
in its image, and none of the $2.2\cdot10^4$ sign and key-placement variants
of this topology---all of them equally stable---is an automorphism.  Such a
scheme evaluates every polynomial only after a shift $x\mapsto x+t$ chosen so
that $\pi\ne0$, which is precisely the generic preconditioning of Motzkin,
Eve, and Knuth; it is not an automorphism of the coefficient space, and it is
the automorphism property that this paper is about.  Within that paradigm the
evidence is uniform: stability was never observed, and
\Cref{prop:unit-pivot-unstable}\ref{it:pivot-sup} shows that a polynomial
automorphism with a descending triangular decoder has unbounded
amplification as soon as one of its rows carries a nonzero correction, so
the majorant bound cannot certify coefficientwise backward stability for
any such decoder.

\paragraph{Open problem.}
Mesztenyi and Witzgall \cite{mesztenyi1967stable} observed in 1967 that
schemes with fewer operations lose significance; Rice
\cite{rice1965conditioning} introduced condition numbers of polynomial forms,
and Miller \cite{miller1975computational} the programme of deriving
complexity lower bounds from stability requirements.  We know of no stability lower bound
for preconditioned polynomial evaluation.  Is there a scheme with
$(1-\epsilon)n$ multiplications and rational preprocessing whose amplification
is $n^{O(1)}$ for all real monic $P$ with $\abs{a_i}\le1$ and $\abs{x}\le1$?
We conjecture not: that $A=n^{O(1)}$ forces $n-O(\log n)$ multiplications,
and that $\floor{n/2}+O(1)$ multiplications force $A\ge\exp(\Omega(n))$.
For polynomial automorphisms the question may be more tractable: we conjecture
that every automorphism with $\floor{n/2}+O(1)$ multiplications has a decoder
of degree exponential in $n$, and already at $n=7$ we know of no stable one.
The measured growth of the Rabin--Winograd amplification up to $n=63$ does
not yet distinguish polynomial from exponential, so even the tradeoff form is
open.
